\documentclass[10pt,a4paper]{article}
\usepackage[margin=20mm,headheight=15pt,headsep=7mm,footskip=10mm]{geometry}
\usepackage{fontspec,amsmath,amssymb,mathtools,microtype}
\setmainfont{lmroman10-regular.otf}[BoldFont=lmroman10-bold.otf,ItalicFont=lmroman10-italic.otf,BoldItalicFont=lmroman10-bolditalic.otf]
\setsansfont{lmsans10-regular.otf}[BoldFont=lmsans10-bold.otf,ItalicFont=lmsans10-oblique.otf,BoldItalicFont=lmsans10-boldoblique.otf]
\setmonofont[Scale=.86]{lmmono10-regular.otf}[ItalicFont=lmmono10-italic.otf]
\usepackage{ragged2e}\usepackage{xcolor,booktabs,longtable,array,graphicx,enumitem,fancyhdr,needspace,titlesec,hyperref,xurl}
\definecolor{AtlasNavy}{HTML}{18334C}\definecolor{AtlasTeal}{HTML}{116D73}\definecolor{AtlasGray}{HTML}{52616D}
\hypersetup{colorlinks=true,linkcolor=AtlasTeal,urlcolor=AtlasTeal,pdfcreator={Complexity Atlas},unicode=true}
\urlstyle{same}\setlength{\parindent}{0pt}\setlength{\parskip}{4.8pt}\setlength{\emergencystretch}{3em}
\renewcommand{\baselinestretch}{1.08}\setcounter{secnumdepth}{0}\setlist{nosep,leftmargin=1.5em}
\titleformat{\section}{\Large\bfseries\sffamily\color{AtlasNavy}\raggedright}{}{0em}{}
\titleformat{\subsection}{\large\bfseries\sffamily\color{AtlasTeal}\raggedright}{}{0em}{}
\titleformat{\subsubsection}{\normalsize\bfseries\color{AtlasNavy}\raggedright}{}{0em}{}
\titlespacing*{\section}{0pt}{14pt}{6pt}\titlespacing*{\subsection}{0pt}{10pt}{5pt}
\pagestyle{fancy}\fancyhf{}\fancyhead[L]{\small\sffamily\color{AtlasGray}COMPLEXITY ATLAS}
\fancyhead[R]{\small\sffamily\color{AtlasGray}ENGLISH EDITION}\fancyfoot[C]{\small\thepage}\renewcommand{\headrulewidth}{.3pt}
\newcommand{\AtlasMath}[1]{\par\begingroup\centering\setbox0=\hbox{$\displaystyle #1$}\ifdim\wd0>1.4\linewidth\typeout{ATLAS WIDE FORMULA: \the\wd0}\fi\ifdim\wd0>\linewidth\resizebox{\linewidth}{!}{\box0}\else\box0\fi\par\endgroup}
\newcommand{\AtlasInlineMath}[1]{\begingroup\setbox0=\hbox{$#1$}\ifdim\wd0>.4\linewidth\endgroup\AtlasMath{#1}\else\endgroup$#1$\fi}
\pdfstringdefDisableCommands{\def\AtlasInlineMath#1{#1}}
\DeclareMathOperator{\polylog}{polylog}\DeclareMathOperator{\dist}{dist}\DeclareMathOperator{\diam}{diam}
\DeclareMathOperator{\rank}{rank}\DeclareMathOperator{\tr}{tr}\DeclareMathOperator{\diag}{diag}
\DeclareMathOperator{\prox}{prox}\DeclareMathOperator{\KL}{KL}\DeclareMathOperator{\Var}{Var}
\providecommand{\eps}{\varepsilon}\providecommand{\R}{\mathbb R}\providecommand{\E}{\mathbb E}
\providecommand{\mathbbm}[1]{\mathbf{#1}}\providecommand{\mathds}[1]{\mathbf{#1}}
\providecommand{\llbracket}{[\![}\providecommand{\rrbracket}{]\!]}
\begin{document}
\raggedright

{\LARGE\bfseries\sffamily\color{AtlasNavy} Global minimization of Hölder-smooth convex functions: initial distance in the full space; exact function values and gradients\par}\vspace{5pt}
\textbf{Setting ID.} \texttt{SO\_\allowbreak{}FO\_\allowbreak{}C05\_\allowbreak{}full\_\allowbreak{}FG}

Single-objective optimization

Find an approximate global minimizer, with accuracy measured by suboptimality of the complete objective.

\textbf{Status.} Polynomial gap present

\textbf{Gap classes present.} Polynomial gap; Other quantified gap; Logarithmic gap

The regional bracket below specifies the quantified gaps and their parameter dependence.

\textbf{TASK tags.} formulation: Global minimization; task family: Single objective

\textbf{MODEL tags.} convexity: Convex; geometry: Full space; geometry: Norm ball; geometry: Norm mismatch allowed; smoothness: Holder smooth

\textbf{ORACLE tags.} access order: First-order interface; feedback quality: Exact; interface family: Fg; operation: Value gradient; public information: Computable access; public information: Hidden objects explicitly declared; public information: Known upper bound; public information: Known exact value; public information: Public inputs explicitly declared

\textbf{ALGORITHMS tags.} adaptivity: Adaptive; computation: Internal arithmetic uncharged; randomization: Deterministic; randomization: Private randomness allowed; restriction: No additional span restriction explicitly imposed

\textbf{OUTPUT tags.} certificate: No additional output certificate required; feasibility: Required; object: Point; postprocessing: Unqueried output explicitly allowed

\textbf{METRIC tags.} criterion: Objective gap

\textbf{GUARANTEE tags.} aggregation: Per instance; mode: Deterministic; mode: Expectation; randomness: Algorithm private seed

\textbf{RESOURCE tags.} budget: Fixed worst-case budget; measure: Oracle calls; precision: Bit precision not counted

\needspace{7\baselineskip}\section{1. Setting specification}
\needspace{4\baselineskip}\subsection{Task}
Find an approximate global minimizer, with accuracy measured by suboptimality of the complete objective.


\AtlasMath{\min_{x\in \mathbb{R}^d}f(x)}


\needspace{4\baselineskip}\subsection{Admissible model}
\AtlasInlineMath{f} is convex on the full space and satisfies \AtlasInlineMath{\|\nabla f(x)-\nabla f(y)\|_{q^\star}\le\beta_\nu\|x-y\|_q^\nu}. A global minimizer satisfies \AtlasInlineMath{\|x^\star-x_0\|_p\le R}. Queries and outputs may lie anywhere in the full space.

Admissible parameter domain:


\AtlasMath{\beta_\nu>0,\quad d\in\mathbb{N}_{\ge1},\quad \varepsilon>0,\quad \nu>0,\quad \nu\le1,\quad R>0,\quad p\ge1,\quad q\ge1}


\needspace{4\baselineskip}\subsection{Oracle interface}
Public numerical inputs: \AtlasInlineMath{d,\ p,\ q,\ \nu,\ \varepsilon,\ x_0}.

Public upper bounds: \AtlasInlineMath{\beta_\nu,\ R}.

The unknown object is the objective \AtlasInlineMath{f}.

Query points lie in \AtlasInlineMath{\mathbb{R}^d}; feedback consists of exact function values and gradients:


\AtlasMath{(f(x),\nabla f(x))}


All responses are consistent with one fixed function instance. Exact values and derivatives agree across repeated queries at the same point.

\needspace{4\baselineskip}\subsection{Output}
Return a point \AtlasInlineMath{\widehat x\in\mathbb R^d}.

The output may be unqueried or constructed from the query history, provided it is feasible. No additional function value or accuracy certificate is required.

\needspace{4\baselineskip}\subsection{Loss or accuracy criterion}

\AtlasMath{\ell(f,x)=f(x)-f^\star}


\AtlasInlineMath{f^\star=\inf_{x\in \mathbb{R}^d}f(x)}. The existence, compactness, or initial-suboptimality assumptions ensure that the optimal value in this comparison is finite.

\needspace{4\baselineskip}\subsection{Resource accounting}
One declared full first-order response costs one call; its vector coordinates are not charged separately.

Cost is bounded uniformly on every execution path. All calls used for initialization, repetition, batching, or candidate verification are counted.

Real arithmetic, precision, and storage applied to already available information are uncharged. Public geometric operations are limited to those explicitly granted by the oracle. Query complexity alone does not imply arithmetic, memory, or communication complexity.

\needspace{4\baselineskip}\subsection{Algorithm views and success guarantees}
\needspace{4\baselineskip}\subsubsection{Deterministic / Pointwise guarantee}
The algorithm uses no private random numbers. Query and stopping decisions may depend adaptively on the observed feedback.

No additional span, zero-respecting, linear-iteration, or memory restriction is imposed.

Interaction is adaptive to the feedback history. Every actual response is charged separately, including responses in parallel or batch queries. Internal real arithmetic, storage, and postprocessing are uncharged.

Select an algorithm using the public parameters, and then require the following for every admissible instance and response rule:


\AtlasMath{\exists A\ \forall I,\mathcal{O}:\ \ell(I,\widehat{x})\le\varepsilon}


The error guarantee holds for every instance; it is not averaged over an input distribution.

\needspace{4\baselineskip}\subsubsection{Randomized / Expected loss}
Private randomization is allowed. Queries, stopping decisions, and postprocessing may adapt to observed feedback.

No additional span, zero-respecting, linear-iteration, or memory restriction is imposed.

Interaction is adaptive to the feedback history. Every actual response is charged separately, including responses in parallel or batch queries. Internal real arithmetic, storage, and postprocessing are uncharged.

Select an algorithm using the public parameters, and then require the following for every admissible instance and response rule:


\AtlasMath{\exists A\ \forall I,\mathcal{O}:\ \mathbb{E}_A[\ell(I,\widehat{x})]\le\varepsilon}


The probability space includes the algorithm's private randomness. The instance and its admissible oracle law are fixed before randomness is averaged.

\needspace{7\baselineskip}\section{2. Notation and parameters}
\AtlasInlineMath{x,\widehat{x}}: Query points and final output; their permitted spaces are specified separately by the oracle and output contracts \textbf{Domain:} \AtlasInlineMath{\widehat{x}\in \mathbb{R}^d}.

\AtlasInlineMath{f}: Unknown objective, with domain and structure specified by the model

\AtlasInlineMath{\nabla f}: Full gradient of the differentiable objective or smooth term with respect to the decision variable

\AtlasInlineMath{x_0}: Public initial point satisfying the stated initial-distance or initial-suboptimality condition

\AtlasInlineMath{f^\star}: Optimal value of the complete task; the oracle need not return it


\AtlasMath{f^\star=\inf_{x\in \mathbb{R}^d}f(x)}


\AtlasInlineMath{x^\star}: A global minimizer whose existence is promised by the model; unique under strong convexity

\AtlasInlineMath{\mathbb{R}^d}: Real vector space with \AtlasInlineMath{d} coordinates

\AtlasInlineMath{\|\cdot\|_2,\langle\cdot,\cdot\rangle}: Euclidean norm and inner product

\AtlasInlineMath{\|\cdot\|_p,\|\cdot\|_q}: Standard finite-dimensional norm with the indicated exponent \textbf{Domain:} Endpoint norm \AtlasInlineMath{\|x\|_\infty=\max_{1\le j\le d}|x_j|}.


\AtlasMath{\|x\|_a=(\sum_{j=1}^d|x_j|^a)^{1/a},\quad1\le a<\infty}


\AtlasInlineMath{d}: Dimension; a public positive integer whose dependence is tracked jointly \textbf{Domain:} \AtlasInlineMath{d\in\mathbb{N}_{\ge1}}.

\AtlasInlineMath{\beta_\nu}: Public upper bound on the gradient Hölder constant \textbf{Domain:} \AtlasInlineMath{\beta_\nu>0}.

\AtlasInlineMath{R}: Upper bound on the distance from the public initial point to a minimizer \textbf{Domain:} \AtlasInlineMath{R>0}.

\AtlasInlineMath{\varepsilon}: Target tolerance for the underlying error defined by the accuracy measure \textbf{Domain:} \AtlasInlineMath{\varepsilon>0}.

\AtlasInlineMath{p}: Public norm exponent specifying the feasible-set or localization geometry \textbf{Domain:} \AtlasInlineMath{1\le p\le\infty}.

\AtlasInlineMath{q}: Public norm exponent specifying the Lipschitz or gradient-regularity geometry \textbf{Domain:} \AtlasInlineMath{1\le q\le\infty}.

\AtlasInlineMath{\nu}: Hölder exponent; public, with fixed or varying status stated in each comparison \textbf{Domain:} \AtlasInlineMath{0<\nu\le1}.

\AtlasInlineMath{c}: Sufficiently small universal positive constant, not specified numerically and potentially different between theorems. \textbf{Domain:} \AtlasInlineMath{0<c<1}.

\AtlasInlineMath{C}: Sufficiently large universal positive constant, not specified numerically and potentially different between theorems. \textbf{Domain:} \AtlasInlineMath{C>1}.

\AtlasInlineMath{\beta}: When \AtlasInlineMath{\nu=1}, write \AtlasInlineMath{\beta=\beta_1}; this does not introduce another independent constant. \textbf{Domain:} \AtlasInlineMath{\beta=\beta_1\quad\text{when }\nu=1}.


\AtlasMath{\beta=\beta_\nu}


\AtlasInlineMath{A_\varepsilon}: Dimensionless accuracy scale for Hölder smoothness


\AtlasMath{A_\varepsilon=\frac{\beta_\nu\,({R}^{1+\nu})}{\varepsilon}}


\AtlasInlineMath{q^\star}: Hölder conjugate exponent of \AtlasInlineMath{q} \textbf{Domain:} \AtlasInlineMath{q^\star\in[1,\infty]}.


\AtlasMath{1/q+1/q^\star=1,\quad1^\star=\infty,\quad\infty^\star=1}


\AtlasInlineMath{K_\nu}: Query scale under Euclidean Hölder smoothness


\AtlasMath{K_\nu={A_\varepsilon}^{\frac{2}{1+3\,\nu}}}


\AtlasInlineMath{T}: Query budget used only in the exact finite-budget result; it is not another model scale. \textbf{Domain:} \AtlasInlineMath{T\in\mathbb{N}_{\ge1}}.

\AtlasInlineMath{S_+}: Error scale after converting the domain and smoothness geometry to Euclidean norms


\AtlasMath{S_+=\beta\,{R}^{2}\,({d}^{2\,(\max\{0,(\frac{1}{q}-0.5)\}+\max\{0,(0.5-\frac{1}{p})\})})}


\AtlasInlineMath{S_+^{(\nu)}}: Euclidean upper-bound error scale for general full-space \AtlasInlineMath{p,q} geometry


\AtlasMath{S_+^{(\nu)}=\beta_\nu\,({R}^{1+\nu})\,({d}^{(1+\nu)\,(\max\{0,(0.5-\frac{1}{p})\}+\max\{0,(\frac{1}{q}-0.5)\})})}


\AtlasInlineMath{S_-^{(\nu)}}: Error scale of Euclidean hard instances embedded in general full-space \AtlasInlineMath{p,q} geometry


\AtlasMath{S_-^{(\nu)}=\beta_\nu\,({R}^{1+\nu})\,({d}^{-1\,(1+\nu)\,(\max\{0,(\frac{1}{p}-0.5)\}+\max\{0,(0.5-\frac{1}{q})\})})}


\AtlasInlineMath{c_{\nu}}: Positive multiplicative constant, not specified numerically, depending only on the following parameters fixed in advance: \AtlasInlineMath{\nu}. \textbf{Domain:} \AtlasInlineMath{c_{\nu}>0}.

\AtlasInlineMath{C_{\nu}}: Positive multiplicative constant, not specified numerically, depending only on the following parameters fixed in advance: \AtlasInlineMath{\nu}. \textbf{Domain:} \AtlasInlineMath{C_{\nu}>0}.

\AtlasInlineMath{C_{p}}: Positive multiplicative constant, not specified numerically, depending only on the following parameters fixed in advance: \AtlasInlineMath{p}. \textbf{Domain:} \AtlasInlineMath{C_{p}>0}.

\AtlasInlineMath{c_{p,\nu}}: Positive multiplicative constant, not specified numerically, depending only on the following parameters fixed in advance: \AtlasInlineMath{p},\AtlasInlineMath{\nu}. \textbf{Domain:} \AtlasInlineMath{c_{p,\nu}>0}.

\AtlasInlineMath{C_{p,\nu}}: Positive multiplicative constant, not specified numerically, depending only on the following parameters fixed in advance: \AtlasInlineMath{p},\AtlasInlineMath{\nu}. \textbf{Domain:} \AtlasInlineMath{C_{p,\nu}>0}.

\AtlasInlineMath{C_{p,q,\nu}}: Positive multiplicative constant, not specified numerically, depending only on the following parameters fixed in advance: \AtlasInlineMath{p},\AtlasInlineMath{q},\AtlasInlineMath{\nu}. \textbf{Domain:} \AtlasInlineMath{C_{p,q,\nu}>0}.

\AtlasInlineMath{\ell}: The underlying loss or error defined by the accuracy measure; its expected, probabilistic, or worst-case aggregation is specified by the success contract

\AtlasInlineMath{I,\mathcal{O},A}: Admissible instance, a response rule consistent with that instance, and an algorithm; the algorithm is selected before the unknown instance

\AtlasInlineMath{O,\Omega,\Theta,\widetilde{O},\widetilde\Omega}: Notation for asymptotic upper and lower bounds, constant-factor matching rates, and explicitly permitted logarithmic factors. Hidden-constant dependence is specified in each regime.

\AtlasInlineMath{\log}: Natural logarithm; changing to another fixed base affects only a constant

\AtlasInlineMath{[a]_+}: Positive part of a real number


\AtlasMath{[a]_+=\max\{a,0\}}


\AtlasInlineMath{\mathbb{E}}: Expectation over the declared randomness after fixing the instance and admissible response law

\needspace{12\baselineskip}\section{3. Complexity bracket}
\needspace{4\baselineskip}\subsection{Compared quantities}
\needspace{4\baselineskip}\subsubsection{Q1 - target accuracy cost}
\textbf{View.} Deterministic / Pointwise guarantee

\AtlasMath{C^\star_{v,\pi}(\varepsilon;\theta)=\inf\{B\ge0:\exists\mathcal A\in\mathfrak A_v,\ C_\pi(\mathcal A,I)\le B\ \text{and}\ G_v(\mathcal A,I;\varepsilon)\ \forall I\}}
\textbf{Definition.} One declared full first-order response costs one call; its vector coordinates are not charged separately.

\textbf{Scope.} declared\_joint\_parameter\_window

\textbf{Optimization direction.} minimize

Here v is the selected algorithm view, G\_v its complete success guarantee, I an admissible fixed instance, theta the remaining parameters, and C\_pi the declared resource projection. C\_pi uses the view's fixed-path or expected-work budget convention, including initialization. No change of guarantee or scalarization is implicit.

\needspace{4\baselineskip}\subsubsection{Q2 - target accuracy cost}
\textbf{View.} Deterministic / Pointwise guarantee

\AtlasMath{C^\star_{v,\pi}(\varepsilon;\theta)=\inf\{B\ge0:\exists\mathcal A\in\mathfrak A_v,\ C_\pi(\mathcal A,I)\le B\ \text{and}\ G_v(\mathcal A,I;\varepsilon)\ \forall I\}}
\textbf{Definition.} One declared full first-order response costs one call; its vector coordinates are not charged separately.

\textbf{Scope.} declared\_joint\_parameter\_window

\textbf{Optimization direction.} minimize

Here v is the selected algorithm view, G\_v its complete success guarantee, I an admissible fixed instance, theta the remaining parameters, and C\_pi the declared resource projection. C\_pi uses the view's fixed-path or expected-work budget convention, including initialization. No change of guarantee or scalarization is implicit.

\needspace{4\baselineskip}\subsubsection{Q3 - target accuracy cost}
\textbf{View.} Deterministic / Pointwise guarantee

\AtlasMath{C^\star_{v,\pi}(\varepsilon;\theta)=\inf\{B\ge0:\exists\mathcal A\in\mathfrak A_v,\ C_\pi(\mathcal A,I)\le B\ \text{and}\ G_v(\mathcal A,I;\varepsilon)\ \forall I\}}
\textbf{Definition.} One declared full first-order response costs one call; its vector coordinates are not charged separately.

\textbf{Optimization direction.} minimize

Here v is the selected algorithm view, G\_v its complete success guarantee, I an admissible fixed instance, theta the remaining parameters, and C\_pi the declared resource projection. C\_pi uses the view's fixed-path or expected-work budget convention, including initialization. No change of guarantee or scalarization is implicit.

\needspace{4\baselineskip}\subsubsection{Q4 - target accuracy cost}
\textbf{View.} Randomized / Expected loss

\AtlasMath{C^\star_{v,\pi}(\varepsilon;\theta)=\inf\{B\ge0:\exists\mathcal A\in\mathfrak A_v,\ C_\pi(\mathcal A,I)\le B\ \text{and}\ G_v(\mathcal A,I;\varepsilon)\ \forall I\}}
\textbf{Definition.} One declared full first-order response costs one call; its vector coordinates are not charged separately.

\textbf{Scope.} declared\_joint\_parameter\_window

\textbf{Optimization direction.} minimize

Here v is the selected algorithm view, G\_v its complete success guarantee, I an admissible fixed instance, theta the remaining parameters, and C\_pi the declared resource projection. C\_pi uses the view's fixed-path or expected-work budget convention, including initialization. No change of guarantee or scalarization is implicit.

\needspace{4\baselineskip}\subsubsection{Q5 - target accuracy cost}
\textbf{View.} Randomized / Expected loss

\AtlasMath{C^\star_{v,\pi}(\varepsilon;\theta)=\inf\{B\ge0:\exists\mathcal A\in\mathfrak A_v,\ C_\pi(\mathcal A,I)\le B\ \text{and}\ G_v(\mathcal A,I;\varepsilon)\ \forall I\}}
\textbf{Definition.} One declared full first-order response costs one call; its vector coordinates are not charged separately.

\textbf{Scope.} declared\_joint\_parameter\_window

\textbf{Optimization direction.} minimize

Here v is the selected algorithm view, G\_v its complete success guarantee, I an admissible fixed instance, theta the remaining parameters, and C\_pi the declared resource projection. C\_pi uses the view's fixed-path or expected-work budget convention, including initialization. No change of guarantee or scalarization is implicit.

\needspace{4\baselineskip}\subsection{R1 - Euclidean Hölder smoothness}
\textbf{Region.} \AtlasInlineMath{p=q=2,\quad0<\nu\le1,\quad\varepsilon\le c_\nu\beta_\nu R^{1+\nu}}

\textbf{Views.} Deterministic / Pointwise guarantee

\textbf{Quantity.} Q1

\textbf{Resource.} One declared full first-order response costs one call; its vector coordinates are not charged separately.

\textbf{Status.} Other quantified gap

\textbf{Lower bound.} B3

\AtlasMath{c_\nu\min\{A_\varepsilon^{2/(1+3\nu)},[d/\log(2+d)]^{1/4}\}}
\textbf{Upper bound.} B1, B2

\AtlasMath{C_\nu\min\{A_\varepsilon^{2/(1+3\nu)},d\log(2+A_\varepsilon)\}}
\textbf{Gap or comparison.} \AtlasInlineMath{\frac{\min\{A_\varepsilon^{2/(1+3\nu)},d\log(2+A_\varepsilon)\}}{\min\{A_\varepsilon^{2/(1+3\nu)},[d/\log(2+d)]^{1/4}\}}}

\textbf{Comparison argument.} The comparison uses the same model, query domain, output requirement, success guarantee, and resource convention. The upper envelope minimizes over applicable algorithms; the lower envelope maximizes over applicable hard instances. A finite sum of lower bounds is used only up to a constant factor. Rate comparison: Both bounds and their validity regions are retained for this contract. An intersection of upper bounds alone does not establish a minimax transition. Ratio: \AtlasInlineMath{\frac{\min\{A_\varepsilon^{2/(1+3\nu)},d\log(2+A_\varepsilon)\}}{\min\{A_\varepsilon^{2/(1+3\nu)},[d/\log(2+d)]^{1/4}\}}}; Permitted fixed-parameter dependence: nu.

\needspace{4\baselineskip}\subsection{R2 - Sufficiently high dimension for the randomized Euclidean lower bound}
\textbf{Region.} \AtlasInlineMath{p=q=2,\quad d\ge C_\nu A_\varepsilon^{8/(1+3\nu)}\log(2+A_\varepsilon),\quad\varepsilon\le c_\nu\beta_\nu R^{1+\nu}}

\textbf{Views.} Deterministic / Pointwise guarantee; Randomized / Expected loss

\textbf{Quantity.} Q1; Q4

\textbf{Resource.} One declared full first-order response costs one call; its vector coordinates are not charged separately.

\textbf{Status.} Optimal

\textbf{Lower bound.} B3, B16

\AtlasMath{c_\nu A_\varepsilon^{2/(1+3\nu)}}
\textbf{Upper bound.} B1, B14

\AtlasMath{C_\nu A_\varepsilon^{2/(1+3\nu)}}
\textbf{Gap or comparison.} \AtlasInlineMath{C_\nu}

\needspace{4\baselineskip}\subsection{R3 - Sufficient randomized dimension for matching geometry with fixed p>=2}
\textbf{Region.} \AtlasInlineMath{2\le p=q<\infty,\quad d\ge C_{p,\nu}A_\varepsilon^{(3p+2)/(1+(p+1)\nu)}\log(2+A_\varepsilon)}

\textbf{Views.} Deterministic / Pointwise guarantee; Randomized / Expected loss

\textbf{Quantity.} Q1; Q4

\textbf{Resource.} One declared full first-order response costs one call; its vector coordinates are not charged separately.

\textbf{Status.} Optimal

\textbf{Lower bound.} B5, B18

\AtlasMath{c_{p,\nu}A_\varepsilon^{p/[1+(p+1)\nu]}}
\textbf{Upper bound.} B4, B2, B17, B15

\AtlasMath{C_{p,\nu}A_\varepsilon^{p/[1+(p+1)\nu]}}
\textbf{Gap or comparison.} \AtlasInlineMath{C_{p,\nu}}

\textbf{Comparison argument.} The comparison uses the same model, query domain, output requirement, success guarantee, and resource convention. The upper envelope minimizes over applicable algorithms; the lower envelope maximizes over applicable hard instances. A finite sum of lower bounds is used only up to a constant factor. Rate comparison: Both bounds and their validity regions are retained for this contract. An intersection of upper bounds alone does not establish a minimax transition. Ratio: \AtlasInlineMath{C_{p,\nu}}; Permitted fixed-parameter dependence: p,nu.

\needspace{4\baselineskip}\subsection{R4 - Sufficient deterministic dimension for matching p<2 geometry}
\textbf{Region.} \AtlasInlineMath{1<p=q<2,\quad d\ge C_{p,\nu}A_\varepsilon^{2/(1+3\nu)}}

\textbf{Views.} Deterministic / Pointwise guarantee

\textbf{Quantity.} Q1

\textbf{Resource.} One declared full first-order response costs one call; its vector coordinates are not charged separately.

\textbf{Status.} Logarithmic gap

\textbf{Lower bound.} B7

\AtlasMath{c_{p,\nu}[A_\varepsilon/\log^\nu(2+d)]^{2/(1+3\nu)}}
\textbf{Upper bound.} B6, B2

\AtlasMath{C_{p,\nu}A_\varepsilon^{2/(1+3\nu)}}
\textbf{Gap or comparison.} \AtlasInlineMath{C_{p,\nu}\log^{2\nu/(1+3\nu)}(2+d)}

\textbf{Comparison argument.} The comparison uses the same model, query domain, output requirement, success guarantee, and resource convention. The upper envelope minimizes over applicable algorithms; the lower envelope maximizes over applicable hard instances. A finite sum of lower bounds is used only up to a constant factor. Rate comparison: Both bounds and their validity regions are retained for this contract. An intersection of upper bounds alone does not establish a minimax transition. Ratio: \AtlasInlineMath{C_{p,\nu}\log^{2\nu/(1+3\nu)}(2+d)}; Permitted fixed-parameter dependence: p,nu.

\needspace{4\baselineskip}\subsection{R5 - Nontrivial full-space norm-embedding regime}
\textbf{Region.} \AtlasInlineMath{0<\varepsilon\le c_\nu S_-^{(\nu)}}

\textbf{Views.} Deterministic / Pointwise guarantee

\textbf{Quantity.} Q2

\textbf{Resource.} One declared full first-order response costs one call; its vector coordinates are not charged separately.

\textbf{Status.} Other quantified gap

\textbf{Lower bound.} B10

\AtlasMath{c_\nu\min\{(S_-^{(\nu)}/\varepsilon)^{2/(1+3\nu)},[d/\log(2+d)]^{1/4}\}}
\textbf{Upper bound.} B9, B2

\AtlasMath{C_\nu[1+\min\{(S_+^{(\nu)}/\varepsilon)^{2/(1+3\nu)},d\log(2+S_+^{(\nu)}/\varepsilon)\}]}
\textbf{Gap or comparison.} 
\AtlasMath{\frac{C_\nu[1+\min\{(S_+^{(\nu)}/\varepsilon)^{2/(1+3\nu)},d\log(2+S_+^{(\nu)}/\varepsilon)\}]}{c_\nu\min\{(S_-^{(\nu)}/\varepsilon)^{2/(1+3\nu)},[d/\log(2+d)]^{1/4}\}}}


\textbf{Comparison argument.} The comparison uses the same model, query domain, output requirement, success guarantee, and resource convention. The upper envelope minimizes over applicable algorithms; the lower envelope maximizes over applicable hard instances. A finite sum of lower bounds is used only up to a constant factor. Rate comparison: This is a conservative norm-embedding baseline valid in every dimension. Matching geometry and nu=1 also admit the stronger specialized comparisons. Ratio: 
\AtlasMath{\frac{C_\nu[1+\min\{(S_+^{(\nu)}/\varepsilon)^{2/(1+3\nu)},d\log(2+S_+^{(\nu)}/\varepsilon)\}]}{c_\nu\min\{(S_-^{(\nu)}/\varepsilon)^{2/(1+3\nu)},[d/\log(2+d)]^{1/4}\}}}
; Permitted fixed-parameter dependence: nu.

\needspace{4\baselineskip}\subsection{R6 - Joint bound in Euclidean geometry with \AtlasInlineMath{\nu=1}}
\textbf{Region.} \AtlasInlineMath{p=q=2,\quad\nu=1,\quad\varepsilon\le c\beta R^2}

\textbf{Views.} Deterministic / Pointwise guarantee

\textbf{Quantity.} Q1

\textbf{Resource.} One declared full first-order response costs one call; its vector coordinates are not charged separately.

\textbf{Status.} Logarithmic gap

\textbf{Lower bound.} B12, B13

\AtlasMath{c\max\{\min\{d,\sqrt{\beta R^2/\varepsilon}\},\log(2+\beta R^2/\varepsilon)\}}
\textbf{Upper bound.} B1, B2

\AtlasMath{C\min\{\sqrt{\beta R^2/\varepsilon},d\log(2+\beta R^2/\varepsilon)\}}
\textbf{Gap or comparison.} \AtlasInlineMath{\frac{C\min\{\sqrt{\beta R^2/\varepsilon},d\log(2+\beta R^2/\varepsilon)\}}{c\max\{\min\{d,\sqrt{\beta R^2/\varepsilon}\},\log(2+\beta R^2/\varepsilon)\}}}

\textbf{Comparison argument.} The comparison uses the same model, query domain, output requirement, success guarantee, and resource convention. The upper envelope minimizes over applicable algorithms; the lower envelope maximizes over applicable hard instances. A finite sum of lower bounds is used only up to a constant factor. Rate comparison: For the slice nu=1, use this complete randomized/deterministic bracket. A generic truncated Hölder lower bound is not the strongest included evidence on this slice. Ratio: \AtlasInlineMath{\frac{C\min\{\sqrt{\beta R^2/\varepsilon},d\log(2+\beta R^2/\varepsilon)\}}{c\max\{\min\{d,\sqrt{\beta R^2/\varepsilon}\},\log(2+\beta R^2/\varepsilon)\}}}

\needspace{4\baselineskip}\subsection{R7 - Sufficient randomized dimension in Euclidean geometry with \AtlasInlineMath{\nu=1}}
\textbf{Region.} \AtlasInlineMath{p=q=2,\nu=1,\quad d\ge C\beta R^2/\varepsilon,\quad\varepsilon\le c\beta R^2}

\textbf{Views.} Deterministic / Pointwise guarantee; Randomized / Expected loss

\textbf{Quantity.} Q1; Q4

\textbf{Resource.} One declared full first-order response costs one call; its vector coordinates are not charged separately.

\textbf{Status.} Optimal

\textbf{Lower bound.} B11, B23

\AtlasMath{c\sqrt{\beta R^2/\varepsilon}}
\textbf{Upper bound.} B1, B14

\AtlasMath{C\sqrt{\beta R^2/\varepsilon}}
\textbf{Gap or comparison.} \AtlasInlineMath{C}

\needspace{4\baselineskip}\subsection{R8 - Sufficient deterministic dimension in Euclidean geometry with \AtlasInlineMath{\nu=1}}
\textbf{Region.} \AtlasInlineMath{p=q=2,\nu=1,\quad d\ge C\sqrt{\beta R^2/\varepsilon},\quad\varepsilon\le c\beta R^2}

\textbf{Views.} Deterministic / Pointwise guarantee

\textbf{Quantity.} Q1

\textbf{Resource.} One declared full first-order response costs one call; its vector coordinates are not charged separately.

\textbf{Status.} Optimal

\textbf{Lower bound.} B12

\AtlasMath{c\sqrt{\beta R^2/\varepsilon}}
\textbf{Upper bound.} B1

\AtlasMath{C\sqrt{\beta R^2/\varepsilon}}
\textbf{Gap or comparison.} \AtlasInlineMath{C}

\needspace{4\baselineskip}\subsection{R9 - Euclidean geometry with \AtlasInlineMath{\nu=1} and \AtlasInlineMath{d=1}}
\textbf{Region.} \AtlasInlineMath{d=1,\quad p=q=2,\quad\nu=1,\quad\varepsilon\le c\beta R^2}

\textbf{Views.} Deterministic / Pointwise guarantee; Randomized / Expected loss

\textbf{Quantity.} Q1; Q4

\textbf{Resource.} One declared full first-order response costs one call; its vector coordinates are not charged separately.

\textbf{Status.} Optimal

\textbf{Lower bound.} B13, B24

\AtlasMath{c\log(2+\beta R^2/\varepsilon)}
\textbf{Upper bound.} B2, B15

\AtlasMath{C\log(2+\beta R^2/\varepsilon)}
\textbf{Gap or comparison.} \AtlasInlineMath{C}

\needspace{4\baselineskip}\subsection{R10 - Euclidean Hölder smoothness}
\textbf{Region.} \AtlasInlineMath{p=q=2,\quad0<\nu\le1,\quad\varepsilon\le c_\nu\beta_\nu R^{1+\nu}}

\textbf{Views.} Randomized / Expected loss

\textbf{Quantity.} Q4

\textbf{Resource.} One declared full first-order response costs one call; its vector coordinates are not charged separately.

\textbf{Status.} Other quantified gap

\textbf{Lower bound.} B16

\AtlasMath{c_\nu\min\{A_\varepsilon^{2/(1+3\nu)},[d/\log(2+d)]^{1/4}\}}
\textbf{Upper bound.} B14, B15

\AtlasMath{C_\nu\min\{A_\varepsilon^{2/(1+3\nu)},d\log(2+A_\varepsilon)\}}
\textbf{Gap or comparison.} \AtlasInlineMath{\frac{\min\{A_\varepsilon^{2/(1+3\nu)},d\log(2+A_\varepsilon)\}}{\min\{A_\varepsilon^{2/(1+3\nu)},[d/\log(2+d)]^{1/4}\}}}

\textbf{Comparison argument.} The comparison uses the same model, query domain, output requirement, success guarantee, and resource convention. The upper envelope minimizes over applicable algorithms; the lower envelope maximizes over applicable hard instances. A finite sum of lower bounds is used only up to a constant factor. Rate comparison: Both bounds and their validity regions are retained for this contract. An intersection of upper bounds alone does not establish a minimax transition. Ratio: \AtlasInlineMath{\frac{\min\{A_\varepsilon^{2/(1+3\nu)},d\log(2+A_\varepsilon)\}}{\min\{A_\varepsilon^{2/(1+3\nu)},[d/\log(2+d)]^{1/4}\}}}; Permitted fixed-parameter dependence: nu.

\needspace{4\baselineskip}\subsection{R11 - Full-space randomized lower-bound window with matching p<2 geometry}
\textbf{Region.} \AtlasInlineMath{1<p=q<2,\quad d\ge C_{p,\nu} A_\varepsilon^2\log(2d)}

\textbf{Views.} Randomized / Expected loss

\textbf{Quantity.} Q4

\textbf{Resource.} One declared full first-order response costs one call; its vector coordinates are not charged separately.

\textbf{Status.} Polynomial gap

\textbf{Lower bound.} B20

\AtlasMath{\frac{c_{p,\nu} A_\varepsilon^{2/(3+2\nu)}}{\log(2+A_\varepsilon)+\nu\log\log(4d)}}
\textbf{Upper bound.} B19, B15

\AtlasMath{C_{p,\nu}\min\{A_\varepsilon^{2/(1+3\nu)},d\log(2+A_\varepsilon)\}}
\textbf{Gap or comparison.} \AtlasInlineMath{C_{p,\nu}A_\varepsilon^{2/(1+3\nu)-2/(3+2\nu)}[\log(2+A_\varepsilon)+\nu\log\log(4d)]}

\textbf{Comparison argument.} The comparison uses the same model, query domain, output requirement, success guarantee, and resource convention. The upper envelope minimizes over applicable algorithms; the lower envelope maximizes over applicable hard instances. A finite sum of lower bounds is used only up to a constant factor. Rate comparison: The accepted full-space initial-distance bracket has a power-law gap. A bounded-query lower bound cannot close it without an appropriate shielding argument. Ratio: \AtlasInlineMath{C_{p,\nu}A_\varepsilon^{2/(1+3\nu)-2/(3+2\nu)}[\log(2+A_\varepsilon)+\nu\log\log(4d)]}; Permitted fixed-parameter dependence: p,nu.

\needspace{4\baselineskip}\subsection{R12 - Nontrivial full-space norm-embedding regime}
\textbf{Region.} \AtlasInlineMath{0<\varepsilon\le c_\nu S_-^{(\nu)}}

\textbf{Views.} Randomized / Expected loss

\textbf{Quantity.} Q5

\textbf{Resource.} One declared full first-order response costs one call; its vector coordinates are not charged separately.

\textbf{Status.} Other quantified gap

\textbf{Lower bound.} B22

\AtlasMath{c_\nu\min\{(S_-^{(\nu)}/\varepsilon)^{2/(1+3\nu)},[d/\log(2+d)]^{1/4}\}}
\textbf{Upper bound.} B21, B15

\AtlasMath{C_\nu[1+\min\{(S_+^{(\nu)}/\varepsilon)^{2/(1+3\nu)},d\log(2+S_+^{(\nu)}/\varepsilon)\}]}
\textbf{Gap or comparison.} 
\AtlasMath{\frac{C_\nu[1+\min\{(S_+^{(\nu)}/\varepsilon)^{2/(1+3\nu)},d\log(2+S_+^{(\nu)}/\varepsilon)\}]}{c_\nu\min\{(S_-^{(\nu)}/\varepsilon)^{2/(1+3\nu)},[d/\log(2+d)]^{1/4}\}}}


\textbf{Comparison argument.} The comparison uses the same model, query domain, output requirement, success guarantee, and resource convention. The upper envelope minimizes over applicable algorithms; the lower envelope maximizes over applicable hard instances. A finite sum of lower bounds is used only up to a constant factor. Rate comparison: This is a conservative norm-embedding baseline valid in every dimension. Matching geometry and nu=1 also admit the stronger specialized comparisons. Ratio: 
\AtlasMath{\frac{C_\nu[1+\min\{(S_+^{(\nu)}/\varepsilon)^{2/(1+3\nu)},d\log(2+S_+^{(\nu)}/\varepsilon)\}]}{c_\nu\min\{(S_-^{(\nu)}/\varepsilon)^{2/(1+3\nu)},[d/\log(2+d)]^{1/4}\}}}
; Permitted fixed-parameter dependence: nu.

\needspace{4\baselineskip}\subsection{R13 - Joint bound in Euclidean geometry with \AtlasInlineMath{\nu=1}}
\textbf{Region.} \AtlasInlineMath{p=q=2,\quad\nu=1,\quad\varepsilon\le c\beta R^2}

\textbf{Views.} Randomized / Expected loss

\textbf{Quantity.} Q4

\textbf{Resource.} One declared full first-order response costs one call; its vector coordinates are not charged separately.

\textbf{Status.} Other quantified gap

\textbf{Lower bound.} B23, B24

\AtlasMath{c\max\{\min\{\sqrt d,\sqrt{\beta R^2/\varepsilon}\},\log(2+\beta R^2/\varepsilon)\}}
\textbf{Upper bound.} B14, B15

\AtlasMath{C\min\{\sqrt{\beta R^2/\varepsilon},d\log(2+\beta R^2/\varepsilon)\}}
\textbf{Gap or comparison.} \AtlasInlineMath{\frac{C\min\{\sqrt{\beta R^2/\varepsilon},d\log(2+\beta R^2/\varepsilon)\}}{c\max\{\min\{\sqrt d,\sqrt{\beta R^2/\varepsilon}\},\log(2+\beta R^2/\varepsilon)\}}}

\textbf{Comparison argument.} The comparison uses the same model, query domain, output requirement, success guarantee, and resource convention. The upper envelope minimizes over applicable algorithms; the lower envelope maximizes over applicable hard instances. A finite sum of lower bounds is used only up to a constant factor. Rate comparison: For the slice nu=1, use this complete randomized/deterministic bracket. A generic truncated Hölder lower bound is not the strongest included evidence on this slice. Ratio: \AtlasInlineMath{\frac{C\min\{\sqrt{\beta R^2/\varepsilon},d\log(2+\beta R^2/\varepsilon)\}}{c\max\{\min\{\sqrt d,\sqrt{\beta R^2/\varepsilon}\},\log(2+\beta R^2/\varepsilon)\}}}

\needspace{7\baselineskip}\section{4. Bounds and supporting evidence}
The following conditions apply to every bound in this profile, in addition to each bound's stated region and algorithm scope.

\textbf{Output mapping.} The output must lie in the declared output space and satisfy its feasibility, support, or atomic-hull restrictions. Permission to return an unqueried point does not remove these restrictions.

\textbf{Resource mapping.} Every primitive response is charged. Sample, component, proximal, and communication costs use the units declared for the target setting. Resource coordinates achieved by one algorithm remain a joint vector; coordinatewise minima from different algorithms are not combined.

\needspace{4\baselineskip}\subsection{B1 - Accelerated gradient method for Hölder-smooth functions}
\textbf{Admission.} Included in the bracket

\textbf{Direction.} upper

\textbf{View.} Deterministic / Pointwise guarantee

\textbf{Witness or lower-bound scope.} Deterministic witness

\AtlasMath{C_\nu(\beta_\nu R^{1+\nu}/\varepsilon)^{2/(1+3\nu)}}
\textbf{Sources.} \href{https://optimization-online.org/wp-content/uploads/2010/12/2865.pdf}{[1]}

\textbf{Conditions.} 
\AtlasMath{p=q=2,\quad0<\nu\le1,\quad\varepsilon\le c_\nu\beta_\nu R^{1+\nu}}


\textbf{Dimension conditions.} \AtlasInlineMath{p=q=2,\quad0<\nu\le1,\quad\varepsilon\le c_\nu\beta_\nu R^{1+\nu}}

\textbf{Oracle correspondence.} Applicable interfaces: exact function values and gradients, exact gradients. Algorithms may use only these interfaces and the auxiliary operations explicitly declared public.

\textbf{Guarantee conversion.} A deterministic error bound for every instance also gives the same expected-error bound. Probability guarantees are used only when an explicit conversion to the target contract is recorded.

\textbf{Source location.} \href{https://optimization-online.org/wp-content/uploads/2010/12/2865.pdf}{[1]} theorem: Section 5.1, Theorem 5, and equation (42); the fixed schedule reads gradients only

\needspace{4\baselineskip}\subsection{B2 - Cutting planes and certificates on the localization ball}
\textbf{Admission.} Included in the bracket

\textbf{Direction.} upper

\textbf{View.} Deterministic / Pointwise guarantee

\textbf{Witness or lower-bound scope.} Deterministic witness

\AtlasMath{C_{p,q,\nu}d\log\left(2+\frac{\beta_\nu R^{1+\nu}d^{(1+\nu)[(1/2-1/p)_++(1/q-1/2)_+]}}{\varepsilon}\right)}
\textbf{Sources.} \href{https://arxiv.org/pdf/2101.09740v2}{[6]} \href{https://arxiv.org/pdf/2609.33360v2}{[7]} \href{https://arxiv.org/pdf/2609.10874}{[4]} \href{https://arxiv.org/pdf/1405.4980}{[8]} \href{https://www2.isye.gatech.edu/~nemirovs/MOR_AccuracyCertificates.pdf}{[9]} \href{https://arxiv.org/pdf/2609.20701}{[10]}

\textbf{Conditions.} 
\AtlasMath{0<\nu\le1,\quad\varepsilon\le c_\nu\beta_\nu R^{1+\nu}}


\textbf{Dimension conditions.} \AtlasInlineMath{0<\nu\le1,\quad\varepsilon\le c_\nu\beta_\nu R^{1+\nu}}

\textbf{Oracle correspondence.} Applicable interfaces: exact function values and gradients, exact gradients. Algorithms may use only these interfaces and the auxiliary operations explicitly declared public.

\textbf{Guarantee conversion.} A deterministic error bound for every instance also gives the same expected-error bound. Probability guarantees are used only when an explicit conversion to the target contract is recorded.

\textbf{Local argument.} A1, A2

\needspace{4\baselineskip}\subsection{B3 - Randomized Hölder-smooth hard instances}
\textbf{Admission.} Included in the bracket

\textbf{Direction.} lower

\textbf{View.} Deterministic / Pointwise guarantee

\textbf{Witness or lower-bound scope.} All randomized adaptive algorithms

\AtlasMath{c_\nu\min\{A_\varepsilon^{2/(1+3\nu)},[d/\log(2+d)]^{1/4}\}}
\textbf{Sources.} \href{https://proceedings.mlr.press/v99/diakonikolas19c/diakonikolas19c.pdf}{[2]}

\textbf{Conditions.} 
\AtlasMath{p=q=2,\quad0<\nu\le1,\quad\varepsilon\le c_\nu\beta_\nu R^{1+\nu}}


\textbf{Dimension conditions.} \AtlasInlineMath{p=q=2,\quad0<\nu\le1,\quad\varepsilon\le c_\nu\beta_\nu R^{1+\nu}}

\textbf{Oracle correspondence.} Applicable interfaces: exact function values and gradients, exact gradients. Lower bounds transfer only through the explicitly recorded simulations from a stronger interface to a weaker one.

\textbf{Guarantee conversion.} The hard-instance quantifiers cover the algorithm class explicitly specified for this bound. Each target setting retains its expected-error or success-probability criterion; a deterministic lower bound is not promoted to a randomized lower bound.

\textbf{Source location.} \href{https://proceedings.mlr.press/v99/diakonikolas19c/diakonikolas19c.pdf}{[2]} theorem: Theorem 12 with \AtlasInlineMath{p=2}, \AtlasInlineMath{K=1}; full-space shielding in equation (3); the stated radius rescaling

\needspace{4\baselineskip}\subsection{B4 - Hölder acceleration for large \AtlasInlineMath{p}}
\textbf{Admission.} Included in the bracket

\textbf{Direction.} upper

\textbf{View.} Deterministic / Pointwise guarantee

\textbf{Witness or lower-bound scope.} Deterministic witness

\AtlasMath{C_{p,\nu}A_\varepsilon^{p/[1+(p+1)\nu]}}
\textbf{Sources.} \href{https://arxiv.org/pdf/1307.5001}{[3]}

\textbf{Conditions.} 
\AtlasMath{2\le p=q<\infty\text{fixed},\quad\varepsilon\le c_{p,\nu}\beta_\nu R^{1+\nu}}


\textbf{Dimension conditions.} \AtlasInlineMath{2\le p=q<\infty\text{fixed},\quad\varepsilon\le c_{p,\nu}\beta_\nu R^{1+\nu}}

\textbf{Oracle correspondence.} Applicable interfaces: exact function values and gradients, exact gradients. Algorithms may use only these interfaces and the auxiliary operations explicitly declared public.

\textbf{Guarantee conversion.} A deterministic error bound for every instance also gives the same expected-error bound. Probability guarantees are used only when an explicit conversion to the target contract is recorded.

\textbf{Source location.} \href{https://arxiv.org/pdf/1307.5001}{[3]} theorem: Section 4.1, the tightness discussion following Corollary 1, and Proposition 3

\needspace{4\baselineskip}\subsection{B5 - Randomized smooth hard instances for large \AtlasInlineMath{p}}
\textbf{Admission.} Included in the bracket

\textbf{Direction.} lower

\textbf{View.} Deterministic / Pointwise guarantee

\textbf{Witness or lower-bound scope.} All randomized adaptive algorithms

\AtlasMath{c_{p,\nu}A_\varepsilon^{p/[1+(p+1)\nu]}}
\textbf{Sources.} \href{https://proceedings.mlr.press/v99/diakonikolas19c/diakonikolas19c.pdf}{[2]}

\textbf{Conditions.} 
\AtlasMath{2\le p=q<\infty,\quad d\ge C_{p,\nu}A_\varepsilon^{(3p+2)/(1+(p+1)\nu)}\log(2+A_\varepsilon)}


\textbf{Dimension conditions.} \AtlasInlineMath{2\le p=q<\infty,\quad d\ge C_{p,\nu}A_\varepsilon^{(3p+2)/(1+(p+1)\nu)}\log(2+A_\varepsilon)}

\textbf{Oracle correspondence.} Applicable interfaces: exact function values and gradients, exact gradients. Lower bounds transfer only through the explicitly recorded simulations from a stronger interface to a weaker one.

\textbf{Guarantee conversion.} The hard-instance quantifiers cover the algorithm class explicitly specified for this bound. Each target setting retains its expected-error or success-probability criterion; a deterministic lower bound is not promoted to a randomized lower bound.

\textbf{Source location.} \href{https://proceedings.mlr.press/v99/diakonikolas19c/diakonikolas19c.pdf}{[2]} theorem: Theorem 12, PDF p. 11; set \AtlasInlineMath{K=1} and a constant success probability, and require the dimension term to be at least the accuracy term

\needspace{4\baselineskip}\subsection{B6 - Universal acceleration for small \AtlasInlineMath{p}}
\textbf{Admission.} Included in the bracket

\textbf{Direction.} upper

\textbf{View.} Deterministic / Pointwise guarantee

\textbf{Witness or lower-bound scope.} Deterministic witness

\AtlasMath{C_{p,\nu}A_\varepsilon^{2/(1+3\nu)}}
\textbf{Sources.} \href{https://optimization-online.org/wp-content/uploads/2010/12/2865.pdf}{[1]} \href{https://optimization-online.org/wp-content/uploads/2013/04/3833.pdf}{[5]}

\textbf{Conditions.} 
\AtlasMath{1<p=q<2\text{fixed},\quad\varepsilon\le c_{p,\nu}\beta_\nu R^{1+\nu}}


\textbf{Dimension conditions.} \AtlasInlineMath{1<p=q<2\text{fixed},\quad\varepsilon\le c_{p,\nu}\beta_\nu R^{1+\nu}}

\textbf{Oracle correspondence.} Applicable interfaces: exact function values and gradients, exact gradients. Algorithms may use only these interfaces and the auxiliary operations explicitly declared public.

\textbf{Guarantee conversion.} A deterministic error bound for every instance also gives the same expected-error bound. Probability guarantees are used only when an explicit conversion to the target contract is recorded.

\textbf{Source location.} \href{https://optimization-online.org/wp-content/uploads/2010/12/2865.pdf}{[1]} theorem: Section 5.1, Theorem 5, and equation (42); the fixed schedule reads gradients only

\textbf{Source location.} \href{https://optimization-online.org/wp-content/uploads/2013/04/3833.pdf}{[5]} theorem: Section 4, specialized to a fixed-parameter mirror map with uniform-convexity order two. The native line search for unknown parameters requires function values.

\needspace{4\baselineskip}\subsection{B7 - Deterministic smooth hard instances for small \AtlasInlineMath{p}}
\textbf{Admission.} Included in the bracket

\textbf{Direction.} lower

\textbf{View.} Deterministic / Pointwise guarantee

\textbf{Witness or lower-bound scope.} All deterministic adaptive algorithms

\AtlasMath{c_{p,\nu}[A_\varepsilon/\log^\nu(2+d)]^{2/(1+3\nu)}}
\textbf{Sources.} \href{https://arxiv.org/pdf/1307.5001}{[3]}

\textbf{Conditions.} 
\AtlasMath{1<p=q<2,\quad d\ge C_{p,\nu}A_\varepsilon^{2/(1+3\nu)}}


\textbf{Dimension conditions.} \AtlasInlineMath{1<p=q<2,\quad d\ge C_{p,\nu}A_\varepsilon^{2/(1+3\nu)}}

\textbf{Oracle correspondence.} Applicable interfaces: exact function values and gradients, exact gradients. Lower bounds transfer only through the explicitly recorded simulations from a stronger interface to a weaker one.

\textbf{Guarantee conversion.} The hard-instance quantifiers cover the algorithm class explicitly specified for this bound. Each target setting retains its expected-error or success-probability criterion; a deterministic lower bound is not promoted to a randomized lower bound.

\textbf{Source location.} \href{https://arxiv.org/pdf/1307.5001}{[3]} theorem: Section 4.1, the tightness discussion following Corollary 1, and Proposition 3

\needspace{4\baselineskip}\subsection{B8 - Randomized smooth hard instances for small \AtlasInlineMath{p}}
\textbf{Admission.} Included in the bracket

\textbf{Direction.} lower

\textbf{View.} Deterministic / Pointwise guarantee

\textbf{Witness or lower-bound scope.} All randomized adaptive algorithms

\AtlasMath{\frac{c_{p,\nu} A_\varepsilon^{2/(3+2\nu)}}{\log(2+A_\varepsilon)+\nu\log\log(4d)}}
\textbf{Sources.} \href{https://proceedings.mlr.press/v99/diakonikolas19c/diakonikolas19c.pdf}{[2]}

\textbf{Conditions.} 
\AtlasMath{1<p=q<2,\quad d\ge C_{p,\nu} A_\varepsilon^2\log(2d)}


\textbf{Dimension conditions.} \AtlasInlineMath{1<p=q<2,\quad d\ge C_{p,\nu} A_\varepsilon^2\log(2d)}

\textbf{Oracle correspondence.} Applicable interfaces: exact function values and gradients, exact gradients. Lower bounds transfer only through the explicitly recorded simulations from a stronger interface to a weaker one.

\textbf{Guarantee conversion.} The hard-instance quantifiers cover the algorithm class explicitly specified for this bound. Each target setting retains its expected-error or success-probability criterion; a deterministic lower bound is not promoted to a randomized lower bound.

\textbf{Source location.} \href{https://proceedings.mlr.press/v99/diakonikolas19c/diakonikolas19c.pdf}{[2]} theorem: Theorem 14, PDF p. 12, and the subsequent sufficient condition \AtlasInlineMath{d\ge C\log(d)\varepsilon^{-2}}

\needspace{4\baselineskip}\subsection{B9 - Hölder acceleration after norm conversion}
\textbf{Admission.} Included in the bracket

\textbf{Direction.} upper

\textbf{View.} Deterministic / Pointwise guarantee

\textbf{Witness or lower-bound scope.} Deterministic witness

\AtlasMath{C_\nu[1+(S_+^{(\nu)}/\varepsilon)^{2/(1+3\nu)}]}
\textbf{Sources.} \href{https://optimization-online.org/wp-content/uploads/2010/12/2865.pdf}{[1]}

\textbf{Conditions.} 
\AtlasMath{0<\varepsilon\le c_\nu S_-^{(\nu)}}


\textbf{Dimension conditions.} \AtlasInlineMath{0<\varepsilon\le c_\nu S_-^{(\nu)}}

\textbf{Oracle correspondence.} Applicable interfaces: exact function values and gradients, exact gradients. Algorithms may use only these interfaces and the auxiliary operations explicitly declared public.

\textbf{Guarantee conversion.} A deterministic error bound for every instance also gives the same expected-error bound. Probability guarantees are used only when an explicit conversion to the target contract is recorded.

\textbf{Source location.} \href{https://optimization-online.org/wp-content/uploads/2010/12/2865.pdf}{[1]} theorem: Section 5.1, Theorem 5, and equation (42); the fixed schedule reads gradients only

\needspace{4\baselineskip}\subsection{B10 - Euclidean randomized hard instances transferred through norm conversion}
\textbf{Admission.} Included in the bracket

\textbf{Direction.} lower

\textbf{View.} Deterministic / Pointwise guarantee

\textbf{Witness or lower-bound scope.} All randomized adaptive algorithms

\AtlasMath{c_\nu\min\{(S_-^{(\nu)}/\varepsilon)^{2/(1+3\nu)},[d/\log(2+d)]^{1/4}\}}
\textbf{Sources.} \href{https://proceedings.mlr.press/v99/diakonikolas19c/diakonikolas19c.pdf}{[2]}

\textbf{Conditions.} 
\AtlasMath{0<\varepsilon\le c_\nu S_-^{(\nu)}}


\textbf{Dimension conditions.} \AtlasInlineMath{0<\varepsilon\le c_\nu S_-^{(\nu)}}

\textbf{Oracle correspondence.} Applicable interfaces: exact function values and gradients, exact gradients. Lower bounds transfer only through the explicitly recorded simulations from a stronger interface to a weaker one.

\textbf{Guarantee conversion.} The hard-instance quantifiers cover the algorithm class explicitly specified for this bound. Each target setting retains its expected-error or success-probability criterion; a deterministic lower bound is not promoted to a randomized lower bound.

\textbf{Source location.} \href{https://proceedings.mlr.press/v99/diakonikolas19c/diakonikolas19c.pdf}{[2]} theorem: Theorem 12 with \AtlasInlineMath{p=2}, \AtlasInlineMath{K=1}; full-space shielding in equation (3); the stated radius rescaling

\needspace{4\baselineskip}\subsection{B11 - Reduction from randomized matrix-vector queries to smooth optimization}
\textbf{Admission.} Included in the bracket

\textbf{Direction.} lower

\textbf{View.} Deterministic / Pointwise guarantee

\textbf{Witness or lower-bound scope.} All randomized adaptive algorithms

\AtlasMath{c\min\{\sqrt d,\sqrt{\beta R^2/\varepsilon}\}}
\textbf{Sources.} \href{https://arxiv.org/pdf/2609.10874}{[4]}

\textbf{Conditions.} 
\AtlasMath{p=q=2,\quad\nu=1,\quad\varepsilon\le c\beta R^2}


\textbf{Dimension conditions.} \AtlasInlineMath{p=q=2,\quad\nu=1,\quad\varepsilon\le c\beta R^2}

\textbf{Oracle correspondence.} Applicable interfaces: exact function values and gradients, exact gradients. Lower bounds transfer only through the explicitly recorded simulations from a stronger interface to a weaker one.

\textbf{Guarantee conversion.} The hard-instance quantifiers cover the algorithm class explicitly specified for this bound. Each target setting retains its expected-error or success-probability criterion; a deterministic lower bound is not promoted to a randomized lower bound.

\textbf{Source location.} \href{https://arxiv.org/pdf/2609.10874}{[4]} theorem: Theorem 1 and equation (7); conversion to the general initial-distance contract

\needspace{4\baselineskip}\subsection{B12 - Deterministic smooth quadratic chain}
\textbf{Admission.} Included in the bracket

\textbf{Direction.} lower

\textbf{View.} Deterministic / Pointwise guarantee

\textbf{Witness or lower-bound scope.} All deterministic adaptive algorithms

\AtlasMath{c\min\{d,\sqrt{\beta R^2/\varepsilon}\}}
\textbf{Sources.} \href{https://arxiv.org/pdf/2101.09740v2}{[6]}

\textbf{Conditions.} 
\AtlasMath{p=q=2,\quad\nu=1,\quad\varepsilon\le c\beta R^2}


\textbf{Dimension conditions.} \AtlasInlineMath{p=q=2,\quad\nu=1,\quad\varepsilon\le c\beta R^2}

\textbf{Oracle correspondence.} Applicable interfaces: exact function values and gradients, exact gradients. Lower bounds transfer only through the explicitly recorded simulations from a stronger interface to a weaker one.

\textbf{Guarantee conversion.} The hard-instance quantifiers cover the algorithm class explicitly specified for this bound. Each target setting retains its expected-error or success-probability criterion; a deterministic lower bound is not promoted to a randomized lower bound.

\textbf{Source location.} \href{https://arxiv.org/pdf/2101.09740v2}{[6]} theorem: Corollary 2, \AtlasInlineMath{d\ge2T+1}, transferred under the same smooth-convex contract

\needspace{4\baselineskip}\subsection{B13 - One-dimensional smooth localization hard instances}
\textbf{Admission.} Included in the bracket

\textbf{Direction.} lower

\textbf{View.} Deterministic / Pointwise guarantee

\textbf{Witness or lower-bound scope.} All randomized adaptive algorithms

\AtlasMath{c\log(2+\beta R^2/\varepsilon)}
\textbf{Sources.} \href{https://arxiv.org/pdf/2101.09740v2}{[6]} \href{https://arxiv.org/pdf/2609.33360v2}{[7]} \href{https://arxiv.org/pdf/2609.10874}{[4]}

\textbf{Conditions.} 
\AtlasMath{p=q=2,\quad\nu=1,\quad\varepsilon\le c\beta R^2}


\textbf{Dimension conditions.} \AtlasInlineMath{p=q=2,\quad\nu=1,\quad\varepsilon\le c\beta R^2}

\textbf{Oracle correspondence.} Applicable interfaces: exact function values and gradients, exact gradients. Lower bounds transfer only through the explicitly recorded simulations from a stronger interface to a weaker one.

\textbf{Guarantee conversion.} The hard-instance quantifiers cover the algorithm class explicitly specified for this bound. Each target setting retains its expected-error or success-probability criterion; a deterministic lower bound is not promoted to a randomized lower bound.

\textbf{Local argument.} A3, A1

\needspace{4\baselineskip}\subsection{B14 - Accelerated gradient method for Hölder-smooth functions}
\textbf{Admission.} Included in the bracket

\textbf{Direction.} upper

\textbf{View.} Randomized / Expected loss

\textbf{Witness or lower-bound scope.} Deterministic witness

\AtlasMath{C_\nu(\beta_\nu R^{1+\nu}/\varepsilon)^{2/(1+3\nu)}}
\textbf{Sources.} \href{https://optimization-online.org/wp-content/uploads/2010/12/2865.pdf}{[1]}

\textbf{Conditions.} 
\AtlasMath{p=q=2,\quad0<\nu\le1,\quad\varepsilon\le c_\nu\beta_\nu R^{1+\nu}}


\textbf{Dimension conditions.} \AtlasInlineMath{p=q=2,\quad0<\nu\le1,\quad\varepsilon\le c_\nu\beta_\nu R^{1+\nu}}

\textbf{Oracle correspondence.} Applicable interfaces: exact function values and gradients, exact gradients. Algorithms may use only these interfaces and the auxiliary operations explicitly declared public.

\textbf{Guarantee conversion.} A deterministic error bound for every instance also gives the same expected-error bound. Probability guarantees are used only when an explicit conversion to the target contract is recorded.

\textbf{Source location.} \href{https://optimization-online.org/wp-content/uploads/2010/12/2865.pdf}{[1]} theorem: Section 5.1, Theorem 5, and equation (42); the fixed schedule reads gradients only

\needspace{4\baselineskip}\subsection{B15 - Cutting planes and certificates on the localization ball}
\textbf{Admission.} Included in the bracket

\textbf{Direction.} upper

\textbf{View.} Randomized / Expected loss

\textbf{Witness or lower-bound scope.} Deterministic witness

\AtlasMath{C_{p,q,\nu}d\log\left(2+\frac{\beta_\nu R^{1+\nu}d^{(1+\nu)[(1/2-1/p)_++(1/q-1/2)_+]}}{\varepsilon}\right)}
\textbf{Sources.} \href{https://arxiv.org/pdf/2101.09740v2}{[6]} \href{https://arxiv.org/pdf/2609.33360v2}{[7]} \href{https://arxiv.org/pdf/2609.10874}{[4]} \href{https://arxiv.org/pdf/1405.4980}{[8]} \href{https://www2.isye.gatech.edu/~nemirovs/MOR_AccuracyCertificates.pdf}{[9]} \href{https://arxiv.org/pdf/2609.20701}{[10]}

\textbf{Conditions.} 
\AtlasMath{0<\nu\le1,\quad\varepsilon\le c_\nu\beta_\nu R^{1+\nu}}


\textbf{Dimension conditions.} \AtlasInlineMath{0<\nu\le1,\quad\varepsilon\le c_\nu\beta_\nu R^{1+\nu}}

\textbf{Oracle correspondence.} Applicable interfaces: exact function values and gradients, exact gradients. Algorithms may use only these interfaces and the auxiliary operations explicitly declared public.

\textbf{Guarantee conversion.} A deterministic error bound for every instance also gives the same expected-error bound. Probability guarantees are used only when an explicit conversion to the target contract is recorded.

\textbf{Local argument.} A1, A2

\needspace{4\baselineskip}\subsection{B16 - Randomized Hölder-smooth hard instances}
\textbf{Admission.} Included in the bracket

\textbf{Direction.} lower

\textbf{View.} Randomized / Expected loss

\textbf{Witness or lower-bound scope.} All randomized adaptive algorithms

\AtlasMath{c_\nu\min\{A_\varepsilon^{2/(1+3\nu)},[d/\log(2+d)]^{1/4}\}}
\textbf{Sources.} \href{https://proceedings.mlr.press/v99/diakonikolas19c/diakonikolas19c.pdf}{[2]}

\textbf{Conditions.} 
\AtlasMath{p=q=2,\quad0<\nu\le1,\quad\varepsilon\le c_\nu\beta_\nu R^{1+\nu}}


\textbf{Dimension conditions.} \AtlasInlineMath{p=q=2,\quad0<\nu\le1,\quad\varepsilon\le c_\nu\beta_\nu R^{1+\nu}}

\textbf{Oracle correspondence.} Applicable interfaces: exact function values and gradients, exact gradients. Lower bounds transfer only through the explicitly recorded simulations from a stronger interface to a weaker one.

\textbf{Guarantee conversion.} The hard-instance quantifiers cover the algorithm class explicitly specified for this bound. Each target setting retains its expected-error or success-probability criterion; a deterministic lower bound is not promoted to a randomized lower bound.

\textbf{Source location.} \href{https://proceedings.mlr.press/v99/diakonikolas19c/diakonikolas19c.pdf}{[2]} theorem: Theorem 12 with \AtlasInlineMath{p=2}, \AtlasInlineMath{K=1}; full-space shielding in equation (3); the stated radius rescaling

\needspace{4\baselineskip}\subsection{B17 - Hölder acceleration for large \AtlasInlineMath{p}}
\textbf{Admission.} Included in the bracket

\textbf{Direction.} upper

\textbf{View.} Randomized / Expected loss

\textbf{Witness or lower-bound scope.} Deterministic witness

\AtlasMath{C_{p,\nu}A_\varepsilon^{p/[1+(p+1)\nu]}}
\textbf{Sources.} \href{https://arxiv.org/pdf/1307.5001}{[3]}

\textbf{Conditions.} 
\AtlasMath{2\le p=q<\infty\text{fixed},\quad\varepsilon\le c_{p,\nu}\beta_\nu R^{1+\nu}}


\textbf{Dimension conditions.} \AtlasInlineMath{2\le p=q<\infty\text{fixed},\quad\varepsilon\le c_{p,\nu}\beta_\nu R^{1+\nu}}

\textbf{Oracle correspondence.} Applicable interfaces: exact function values and gradients, exact gradients. Algorithms may use only these interfaces and the auxiliary operations explicitly declared public.

\textbf{Guarantee conversion.} A deterministic error bound for every instance also gives the same expected-error bound. Probability guarantees are used only when an explicit conversion to the target contract is recorded.

\textbf{Source location.} \href{https://arxiv.org/pdf/1307.5001}{[3]} theorem: Section 4.1, the tightness discussion following Corollary 1, and Proposition 3

\needspace{4\baselineskip}\subsection{B18 - Randomized smooth hard instances for large \AtlasInlineMath{p}}
\textbf{Admission.} Included in the bracket

\textbf{Direction.} lower

\textbf{View.} Randomized / Expected loss

\textbf{Witness or lower-bound scope.} All randomized adaptive algorithms

\AtlasMath{c_{p,\nu}A_\varepsilon^{p/[1+(p+1)\nu]}}
\textbf{Sources.} \href{https://proceedings.mlr.press/v99/diakonikolas19c/diakonikolas19c.pdf}{[2]}

\textbf{Conditions.} 
\AtlasMath{2\le p=q<\infty,\quad d\ge C_{p,\nu}A_\varepsilon^{(3p+2)/(1+(p+1)\nu)}\log(2+A_\varepsilon)}


\textbf{Dimension conditions.} \AtlasInlineMath{2\le p=q<\infty,\quad d\ge C_{p,\nu}A_\varepsilon^{(3p+2)/(1+(p+1)\nu)}\log(2+A_\varepsilon)}

\textbf{Oracle correspondence.} Applicable interfaces: exact function values and gradients, exact gradients. Lower bounds transfer only through the explicitly recorded simulations from a stronger interface to a weaker one.

\textbf{Guarantee conversion.} The hard-instance quantifiers cover the algorithm class explicitly specified for this bound. Each target setting retains its expected-error or success-probability criterion; a deterministic lower bound is not promoted to a randomized lower bound.

\textbf{Source location.} \href{https://proceedings.mlr.press/v99/diakonikolas19c/diakonikolas19c.pdf}{[2]} theorem: Theorem 12, PDF p. 11; set \AtlasInlineMath{K=1} and a constant success probability, and require the dimension term to be at least the accuracy term

\needspace{4\baselineskip}\subsection{B19 - Universal acceleration for small \AtlasInlineMath{p}}
\textbf{Admission.} Included in the bracket

\textbf{Direction.} upper

\textbf{View.} Randomized / Expected loss

\textbf{Witness or lower-bound scope.} Deterministic witness

\AtlasMath{C_{p,\nu}A_\varepsilon^{2/(1+3\nu)}}
\textbf{Sources.} \href{https://optimization-online.org/wp-content/uploads/2010/12/2865.pdf}{[1]} \href{https://optimization-online.org/wp-content/uploads/2013/04/3833.pdf}{[5]}

\textbf{Conditions.} 
\AtlasMath{1<p=q<2\text{fixed},\quad\varepsilon\le c_{p,\nu}\beta_\nu R^{1+\nu}}


\textbf{Dimension conditions.} \AtlasInlineMath{1<p=q<2\text{fixed},\quad\varepsilon\le c_{p,\nu}\beta_\nu R^{1+\nu}}

\textbf{Oracle correspondence.} Applicable interfaces: exact function values and gradients, exact gradients. Algorithms may use only these interfaces and the auxiliary operations explicitly declared public.

\textbf{Guarantee conversion.} A deterministic error bound for every instance also gives the same expected-error bound. Probability guarantees are used only when an explicit conversion to the target contract is recorded.

\textbf{Source location.} \href{https://optimization-online.org/wp-content/uploads/2010/12/2865.pdf}{[1]} theorem: Section 5.1, Theorem 5, and equation (42); the fixed schedule reads gradients only

\textbf{Source location.} \href{https://optimization-online.org/wp-content/uploads/2013/04/3833.pdf}{[5]} theorem: Section 4, specialized to a fixed-parameter mirror map with uniform-convexity order two. The native line search for unknown parameters requires function values.

\needspace{4\baselineskip}\subsection{B20 - Randomized smooth hard instances for small \AtlasInlineMath{p}}
\textbf{Admission.} Included in the bracket

\textbf{Direction.} lower

\textbf{View.} Randomized / Expected loss

\textbf{Witness or lower-bound scope.} All randomized adaptive algorithms

\AtlasMath{\frac{c_{p,\nu} A_\varepsilon^{2/(3+2\nu)}}{\log(2+A_\varepsilon)+\nu\log\log(4d)}}
\textbf{Sources.} \href{https://proceedings.mlr.press/v99/diakonikolas19c/diakonikolas19c.pdf}{[2]}

\textbf{Conditions.} 
\AtlasMath{1<p=q<2,\quad d\ge C_{p,\nu} A_\varepsilon^2\log(2d)}


\textbf{Dimension conditions.} \AtlasInlineMath{1<p=q<2,\quad d\ge C_{p,\nu} A_\varepsilon^2\log(2d)}

\textbf{Oracle correspondence.} Applicable interfaces: exact function values and gradients, exact gradients. Lower bounds transfer only through the explicitly recorded simulations from a stronger interface to a weaker one.

\textbf{Guarantee conversion.} The hard-instance quantifiers cover the algorithm class explicitly specified for this bound. Each target setting retains its expected-error or success-probability criterion; a deterministic lower bound is not promoted to a randomized lower bound.

\textbf{Source location.} \href{https://proceedings.mlr.press/v99/diakonikolas19c/diakonikolas19c.pdf}{[2]} theorem: Theorem 14, PDF p. 12, and the subsequent sufficient condition \AtlasInlineMath{d\ge C\log(d)\varepsilon^{-2}}

\needspace{4\baselineskip}\subsection{B21 - Hölder acceleration after norm conversion}
\textbf{Admission.} Included in the bracket

\textbf{Direction.} upper

\textbf{View.} Randomized / Expected loss

\textbf{Witness or lower-bound scope.} Deterministic witness

\AtlasMath{C_\nu[1+(S_+^{(\nu)}/\varepsilon)^{2/(1+3\nu)}]}
\textbf{Sources.} \href{https://optimization-online.org/wp-content/uploads/2010/12/2865.pdf}{[1]}

\textbf{Conditions.} 
\AtlasMath{0<\varepsilon\le c_\nu S_-^{(\nu)}}


\textbf{Dimension conditions.} \AtlasInlineMath{0<\varepsilon\le c_\nu S_-^{(\nu)}}

\textbf{Oracle correspondence.} Applicable interfaces: exact function values and gradients, exact gradients. Algorithms may use only these interfaces and the auxiliary operations explicitly declared public.

\textbf{Guarantee conversion.} A deterministic error bound for every instance also gives the same expected-error bound. Probability guarantees are used only when an explicit conversion to the target contract is recorded.

\textbf{Source location.} \href{https://optimization-online.org/wp-content/uploads/2010/12/2865.pdf}{[1]} theorem: Section 5.1, Theorem 5, and equation (42); the fixed schedule reads gradients only

\needspace{4\baselineskip}\subsection{B22 - Euclidean randomized hard instances transferred through norm conversion}
\textbf{Admission.} Included in the bracket

\textbf{Direction.} lower

\textbf{View.} Randomized / Expected loss

\textbf{Witness or lower-bound scope.} All randomized adaptive algorithms

\AtlasMath{c_\nu\min\{(S_-^{(\nu)}/\varepsilon)^{2/(1+3\nu)},[d/\log(2+d)]^{1/4}\}}
\textbf{Sources.} \href{https://proceedings.mlr.press/v99/diakonikolas19c/diakonikolas19c.pdf}{[2]}

\textbf{Conditions.} 
\AtlasMath{0<\varepsilon\le c_\nu S_-^{(\nu)}}


\textbf{Dimension conditions.} \AtlasInlineMath{0<\varepsilon\le c_\nu S_-^{(\nu)}}

\textbf{Oracle correspondence.} Applicable interfaces: exact function values and gradients, exact gradients. Lower bounds transfer only through the explicitly recorded simulations from a stronger interface to a weaker one.

\textbf{Guarantee conversion.} The hard-instance quantifiers cover the algorithm class explicitly specified for this bound. Each target setting retains its expected-error or success-probability criterion; a deterministic lower bound is not promoted to a randomized lower bound.

\textbf{Source location.} \href{https://proceedings.mlr.press/v99/diakonikolas19c/diakonikolas19c.pdf}{[2]} theorem: Theorem 12 with \AtlasInlineMath{p=2}, \AtlasInlineMath{K=1}; full-space shielding in equation (3); the stated radius rescaling

\needspace{4\baselineskip}\subsection{B23 - Reduction from randomized matrix-vector queries to smooth optimization}
\textbf{Admission.} Included in the bracket

\textbf{Direction.} lower

\textbf{View.} Randomized / Expected loss

\textbf{Witness or lower-bound scope.} All randomized adaptive algorithms

\AtlasMath{c\min\{\sqrt d,\sqrt{\beta R^2/\varepsilon}\}}
\textbf{Sources.} \href{https://arxiv.org/pdf/2609.10874}{[4]}

\textbf{Conditions.} 
\AtlasMath{p=q=2,\quad\nu=1,\quad\varepsilon\le c\beta R^2}


\textbf{Dimension conditions.} \AtlasInlineMath{p=q=2,\quad\nu=1,\quad\varepsilon\le c\beta R^2}

\textbf{Oracle correspondence.} Applicable interfaces: exact function values and gradients, exact gradients. Lower bounds transfer only through the explicitly recorded simulations from a stronger interface to a weaker one.

\textbf{Guarantee conversion.} The hard-instance quantifiers cover the algorithm class explicitly specified for this bound. Each target setting retains its expected-error or success-probability criterion; a deterministic lower bound is not promoted to a randomized lower bound.

\textbf{Source location.} \href{https://arxiv.org/pdf/2609.10874}{[4]} theorem: Theorem 1 and equation (7); conversion to the general initial-distance contract

\needspace{4\baselineskip}\subsection{B24 - One-dimensional smooth localization hard instances}
\textbf{Admission.} Included in the bracket

\textbf{Direction.} lower

\textbf{View.} Randomized / Expected loss

\textbf{Witness or lower-bound scope.} All randomized adaptive algorithms

\AtlasMath{c\log(2+\beta R^2/\varepsilon)}
\textbf{Sources.} \href{https://arxiv.org/pdf/2101.09740v2}{[6]} \href{https://arxiv.org/pdf/2609.33360v2}{[7]} \href{https://arxiv.org/pdf/2609.10874}{[4]}

\textbf{Conditions.} 
\AtlasMath{p=q=2,\quad\nu=1,\quad\varepsilon\le c\beta R^2}


\textbf{Dimension conditions.} \AtlasInlineMath{p=q=2,\quad\nu=1,\quad\varepsilon\le c\beta R^2}

\textbf{Oracle correspondence.} Applicable interfaces: exact function values and gradients, exact gradients. Lower bounds transfer only through the explicitly recorded simulations from a stronger interface to a weaker one.

\textbf{Guarantee conversion.} The hard-instance quantifiers cover the algorithm class explicitly specified for this bound. Each target setting retains its expected-error or success-probability criterion; a deterministic lower bound is not promoted to a randomized lower bound.

\textbf{Local argument.} A3, A1

\needspace{7\baselineskip}\section{5. Local arguments}
\needspace{4\baselineskip}\subsection{A1 - Smooth localization and gradient-norm conversion}
\textbf{Scope.} Gradient access and a public ball localizing a global minimizer; public regularization when only initial gap is given.

\textbf{Algorithm scope.} Deterministic upper witnesses; imported lower theorems retain their own deterministic or randomized scopes.

If a global minimizer lies within distance \AtlasInlineMath{R} of \AtlasInlineMath{x_0}, query on \AtlasInlineMath{B(x_0,R)}. Since \AtlasInlineMath{\nabla f(x^\star)=0}, gradients there have norm at most \AtlasInlineMath{2\beta R}. The centroid certificate gives \AtlasInlineMath{O(d\log(2+\beta R^2/\epsilon))} gradient calls. Smoothness gives \AtlasInlineMath{f(x_0)-f^\star\le\beta R^2/2}, including the zero-query endpoint.

With global \AtlasInlineMath{\mu}-strong convexity and initial gap \AtlasInlineMath{\Delta}, use radius \AtlasInlineMath{\sqrt{2\Delta/\mu}}, obtaining \AtlasInlineMath{O(d\log(2+\kappa\Delta/\epsilon))}. At \AtlasInlineMath{\mu=\beta}, equality in strong monotonicity and Lipschitz continuity forces \AtlasInlineMath{\nabla f(x)-\nabla f(y)=\beta(x-y)}. One gradient identifies the minimizer of this known-curvature quadratic plus unknown affine term. Opposite translations exclude zero calls at nontrivial accuracy.

Under global Hölder smoothness in \AtlasInlineMath{\ell_q}, convert to Euclidean quantities
\AtlasInlineMath{\beta_E=\beta_\nu d^{(1+\nu)(1/q-1/2)_+}} and
\AtlasInlineMath{R_E=Rd^{(1/2-1/p)_+}}. Gradients on the localization ball are at most \AtlasInlineMath{\beta_E(2R_E)^\nu}, giving
\AtlasInlineMath{C_\nu d\log(2+\beta_E R_E^{1+\nu}/\epsilon)}.
A boundary minimizer of an arbitrary constrained problem does not justify the zero-gradient premise. With two public norm-distance bounds, their intersection lies in the Euclidean ball of radius \AtlasInlineMath{R_{2,+}}, while the Euclidean ball of radius \AtlasInlineMath{R_{2,-}} satisfies both promises. The former gives a certificate upper bound; embedding the one-dimensional hard family in the latter gives the logarithmic lower bound.

For gradient-norm accuracy, use \AtlasInlineMath{\|\nabla f(x)\|^2\le2\beta(f(x)-f^\star)}. Under initial distance, target gap \AtlasInlineMath{\epsilon^2/(2\beta)}. Under initial gap alone, instead set

\AtlasMath{F_\lambda=f+\lambda\|x-x_0\|^2/2,\qquad\lambda=\epsilon^2/(32\Delta).}

Its minimizer lies within \AtlasInlineMath{8\Delta/\epsilon}. A certificate at gap \AtlasInlineMath{\epsilon^2/(128\beta)} gives \AtlasInlineMath{\|\nabla F_\lambda(\widehat x)\|\le\epsilon/4}, while the public regularizer contributes at most \AtlasInlineMath{\epsilon/4}. Each regularized gradient costs one original gradient, and the budget is \AtlasInlineMath{O(d\log(2+\beta\Delta/\epsilon^2))}.

For lower transfers, strong convexity gives \AtlasInlineMath{f-f^\star\le\|\nabla f\|^2/(2\mu)}. Choose the hard subclass modulus proportional to \AtlasInlineMath{\epsilon/R}, or \AtlasInlineMath{\epsilon^2/\Delta}, for the two respective initializations. Gradient accuracy forces constant relative gap reduction. Apply the cited deterministic chain or randomized matrix-query theorem in its original dimension window; truncate only to constructions that fit. This preserves initialization, guarantee, and algorithm scope.

For an already strongly convex model, the gap threshold required for the lower transfer is \AtlasInlineMath{\epsilon^2/(2\mu)}, whereas the sufficient gap for the upper transfer is \AtlasInlineMath{\epsilon^2/(2\beta)}. Retain these different logarithm arguments rather than silently canceling their condition-number shift. For \AtlasInlineMath{1<\kappa\le4}, gradient descent and the nested one-dimensional construction yield the stated comparable near-unit denominators. At \AtlasInlineMath{\kappa=1}, the known-curvature quadratic is solved in one call.

For matching \AtlasInlineMath{\ell_p} Hölder geometry, the descent model along a dual steepest-descent direction gives
\AtlasInlineMath{f(x)-f^\star\ge c_\nu\beta_\nu^{-1/\nu}\|\nabla f(x)\|_{p^\star}^{(1+\nu)/\nu}}.
Targeting gap \AtlasInlineMath{c_\nu\beta_\nu^{-1/\nu}\epsilon^{(1+\nu)/\nu}} in the localization certificate therefore gives
\AtlasInlineMath{C_\nu d\log(2+\beta_\nu R^\nu/\epsilon)}.

\textbf{Limits of the argument.} An arbitrary constrained minimizer may be on the boundary and need not have zero gradient. Fixed-dimensional lower transfers retain theorem dimension windows.

\textbf{Sources.} \href{https://arxiv.org/pdf/2101.09740v2}{[6]} \href{https://arxiv.org/pdf/2609.33360v2}{[7]} \href{https://arxiv.org/pdf/2609.10874}{[4]}

\needspace{4\baselineskip}\subsection{A2 - Subgradient certificates and centroid cuts}
\textbf{Scope.} One full subgradient is sufficient. Public compact convex geometry and exact internal geometric computation are available.

\textbf{Algorithm scope.} Deterministic upper-bound witness; valid as an upper bound for randomized views.

Let \AtlasInlineMath{Q\subset\mathbb R^d} be public, compact, convex, and full dimensional, with diameter at most \AtlasInlineMath{D} in the stated norm. Suppose all returned subgradients have dual norm at most \AtlasInlineMath{L}. Query \AtlasInlineMath{g_i\in\partial f(x_i)} at the centroid of the current body and retain \AtlasInlineMath{\langle g_i,x_i-x\rangle\ge0}. The centroid halfspace theorem bounds the retained volume by \AtlasInlineMath{(1-e^{-1})^T\operatorname{vol}(Q)}.

The computable certificate is

\AtlasMath{\Gamma_T=\min_{\lambda\in\Delta_T}\max_{x\in Q}\sum_i\lambda_i\langle g_i,x_i-x\rangle
=\max_{x\in Q}\min_i\langle g_i,x_i-x\rangle.}

The equality is convex minimax. Subgradient inequalities and convexity imply
\AtlasInlineMath{f(\sum_i\lambda_ix_i)-\min_Qf\le\Gamma_T}.
If \AtlasInlineMath{\Gamma_T>\epsilon}, choose \AtlasInlineMath{z\in Q} with all cut margins greater than \AtlasInlineMath{\epsilon}. The copy \AtlasInlineMath{(1-a)z+aQ}, where \AtlasInlineMath{a=\epsilon/[2(\epsilon+LD)]}, survives every cut. Its volume is \AtlasInlineMath{a^d\operatorname{vol}(Q)}, contradicting centroid shrinkage once \AtlasInlineMath{T>C d\log(2+LD/\epsilon)}. If the body loses volume earlier, the same positive-volume argument already certifies the error. No function values or inscribed-ball estimate are required.

For \AtlasInlineMath{Q=RB_p^d} and \AtlasInlineMath{\ell_q}-Lipschitz objectives, use \AtlasInlineMath{D=2Rd^{(1/q-1/p)_+}}. Returning the center has error at most \AtlasInlineMath{LRd^{(1/q-1/p)_+}}, giving the coarse zero-query regime. Exact centroid, minimax, and public geometric calculations are uncharged internal operations.

For a deterministic selector with total movement at most \AtlasInlineMath{W(T)} on nested convex sets, instead impose deep cuts of margin \AtlasInlineMath{2LW(T)/T}. If all \AtlasInlineMath{T} sets are nonempty, every move has length at least \AtlasInlineMath{2W(T)/T}, contradicting the movement bound. At the first empty intersection, minimax gives error at most \AtlasInlineMath{2LW(T)/T}. Substitution of the cited stable-movement theorem yields the recorded nondual upper bound. Its deterministic information scope is distinct from the scope of randomized implementations.

\textbf{Limits of the argument.} The stable-selector branch imports its movement theorem and does not claim efficient deterministic arithmetic time.

\textbf{Sources.} \href{https://arxiv.org/pdf/1405.4980}{[8]} \href{https://www2.isye.gatech.edu/~nemirovs/MOR_AccuracyCertificates.pdf}{[9]} \href{https://arxiv.org/pdf/2609.20701}{[10]}

\needspace{4\baselineskip}\subsection{A3 - A one-dimensional lower bound with exact values and gradients}
\textbf{Scope.} Exact values and gradients; unbounded real-line queries, with higher-dimensional embeddings.

\textbf{Algorithm scope.} Lower bounds against arbitrary randomized adaptive algorithms, including unqueried output.

Fix \AtlasInlineMath{1<\kappa=\beta/\mu}. Choose a fixed smooth \AtlasInlineMath{\psi}, equal to \AtlasInlineMath{-x} on \AtlasInlineMath{[-1/2,1/2]} and zero outside \AtlasInlineMath{[-1,1]}, and let \AtlasInlineMath{C_\psi=\max\{1,\|\psi''\|_\infty\}}. Set

\AtlasMath{\beta_0=(\beta+\mu)/2,\quad a={\kappa-1\over8C_\psi(\kappa+1)},\quad
\eta=a/4,\quad \rho_0={(1-\eta)\over a}\sqrt{\Delta/\beta}.}

For independent uniform signs \AtlasInlineMath{b_i}, define \AtlasInlineMath{u_0=0},
\AtlasInlineMath{u_i=u_{i-1}+b_i a\rho_{i-1}}, \AtlasInlineMath{\rho_i=\eta\rho_{i-1}}, and

\AtlasMath{f_b(x)=\beta_0x^2/2+
\beta_0\sum_{i=1}^N b_i a\rho_{i-1}^2
\psi((x-u_{i-1})/\rho_{i-1}).}

The two child intervals \AtlasInlineMath{I_i=[u_i-\rho_i,u_i+\rho_i]} are disjoint and lie inside the central half of the parent. Second derivatives of successive bumps have disjoint shell supports. Thus \AtlasInlineMath{\mu\le f_b''\le\beta}, and \AtlasInlineMath{f_b''\ge7\beta/16}. Within \AtlasInlineMath{I_i}, the partial sum through level \AtlasInlineMath{i} is a quadratic centered at \AtlasInlineMath{u_i}. The final minimizer is \AtlasInlineMath{u_N}, with \AtlasInlineMath{|u_N|\le\sqrt{\Delta/\beta}} and initial gap at most \AtlasInlineMath{\Delta/2}.

A stronger oracle reveals the next bit when the query lies in the current unresolved interval, continuing only while it lies in the selected child. Unrevealed bumps vanish near the query, so the revealed prefix computes the exact value and derivative. Conditional on history, a fixed point belongs to at most one child; one query reveals at most \AtlasInlineMath{\sum_{j\ge1}2^{-(j-1)}=2} new bits in expectation.

Treat the final output as one more such query. If \AtlasInlineMath{|\widehat x-u_N|<\rho_N}, all \AtlasInlineMath{N} bits are revealed. Its probability is at most \AtlasInlineMath{2(T+1)/N}. Otherwise the loss is at least
\AtlasInlineMath{c\beta\rho_N^2=c'\Delta\eta^{2(N-1)}}.
Choose \AtlasInlineMath{N} at the accuracy scale and use opposite translations to exclude zero calls. For expected error and constant-probability accuracy,

\AtlasMath{T\ge c\left[1+{\log(\Delta/\epsilon)\over
\log(C\kappa/(\kappa-1))}\right].}

Queries may be arbitrary, randomized, adaptive, and unbounded.

Fixing \AtlasInlineMath{\kappa=2}, \AtlasInlineMath{\Delta=c\beta R^2}, gives the initial-distance logarithmic bound. The gradient version uses
\AtlasInlineMath{|f_b'(x)|\ge(7\beta/16)|x-u_N|}.
Add a public quadratic to embed in higher dimension. For composite minimization, apply the construction on total curvature interval \AtlasInlineMath{[\mu,\mu+\beta]}, with public \AtlasInlineMath{h(x)=\mu x^2/2}; the unknown remainder is convex and \AtlasInlineMath{\beta}-smooth, and every proximal response for \AtlasInlineMath{h} is independent of the hidden bits.

\textbf{Limits of the argument.} Curvature, initialization, and nontrivial-accuracy restrictions remain explicit; the composite transfer lower-bounds unknown-gradient calls.

\needspace{7\baselineskip}\section{References}
[1] Olivier Devolder; François Glineur; Yurii Nesterov. First-order methods of smooth convex optimization with inexact oracle. 2014; public preprint in 2011.

\url{https://optimization-online.org/wp-content/uploads/2010/12/2865.pdf}

[2] Jelena Diakonikolas; Cristóbal Guzmán. Lower Bounds for Parallel and Randomized Convex Optimization. 2019.

\url{https://proceedings.mlr.press/v99/diakonikolas19c/diakonikolas19c.pdf}

[3] Cristóbal Guzmán; Arkadi Nemirovski. On Lower Complexity Bounds for Large-Scale Smooth Convex Optimization. 2015; arXiv2013.

\url{https://arxiv.org/pdf/1307.5001}

[4] Ohad Shamir. A Simple Complexity Lower Bound for Solving Ax=b. 2026-09-09; preprintv1.

\url{https://arxiv.org/pdf/2609.10874}

[5] Yurii Nesterov. Universal Gradient Methods for Convex Optimization Problems. 2015; CORE preprint in 2013.

\url{https://optimization-online.org/wp-content/uploads/2013/04/3833.pdf}

[6] Yoel Drori; Adrien Taylor. On the oracle complexity of smooth strongly convex minimization. 2021; arXiv:2101.09740v2.

\url{https://arxiv.org/pdf/2101.09740v2}

[7] Yoel Drori. Exact oracle complexity for function-value suboptimality under weighted initial conditions. 2026-10-16; preprintv2.

\url{https://arxiv.org/pdf/2609.33360v2}

[8] Sébastien Bubeck. Convex Optimization: Algorithms and Complexity. 2015.

\url{https://arxiv.org/pdf/1405.4980}

[9] Arkadi Nemirovski; Shmuel Onn; Uriel G. Rothblum. Accuracy Certificates for Computational Problems with Convex Structure. 2010.

\url{https://www2.isye.gatech.edu/~nemirovs/MOR_AccuracyCertificates.pdf}

[10] David Martínez-Rubio; Cristóbal Guzmán. Stable Movement for Nondual Lipschitz Convex Optimization: Efficiency and Nearly Optimal Oracle Rates. 2026-09-17; preprintv1.

\url{https://arxiv.org/pdf/2609.20701}

\end{document}
